% \chapterfour

\section{Техническая документация и руководство пользователей}

\subsection{Роли пользователей системы}

В системе есть два типа пользователей, и каждый взаимодействует с ней по-своему.

\textbf{Администратор системы} — это технический специалист, который разворачивает и поддерживает работу системы. Ему нужно из корневой папки проекта выполнить команду \texttt{docker compose up -d}, чтобы всё запустилось. Также администратор настраивает переменные окружения в файлах \texttt{.env} для каждого сервиса: там указываются адреса сервисов, параметры подключения к базе данных, какая LLM-модель используется, размер чанков и так далее. Он следит за состоянием сервисов и очередью задач, управляет объектным хранилищем MinIO и настраивает резервное копирование PostgreSQL.

\textbf{Конечный пользователь} — это аналитик или любой другой сотрудник организации, который работает с системой через веб-интерфейс и не обязан разбираться в технических деталях. Его сценарии работы описаны ниже.

\subsection{Руководство конечного пользователя}

\subsubsection{Загрузка документов в базу знаний}

Чтобы загрузить документ, откройте веб-интерфейс в браузере, перейдите в раздел \textbf{«Документы»} и нажмите кнопку \textbf{«Загрузить документ»}. Система принимает файлы в форматах PDF, DOCX и TXT. Сразу после загрузки система подтверждает, что файл получен, и показывает статус обработки: сначала «В очереди», потом «Обрабатывается», и наконец «Готов». Когда статус сменится на «Готов», документ проиндексирован и его можно искать.

\begin{quote}
\textbf{Важно:} если документ большой (от 100 страниц), обработка может занять несколько минут. Вы можете спокойно продолжать работу с системой — индексация идёт в фоновом режиме.
\end{quote}

\subsubsection{Работа с чатом}

Конечный пользователь заводит чат, и общается с системой на естественном языке. Например, он может спросить: \textit{"Какие билеты есть на 23 февраля ?"}, то система автоматически выполнит семантический поиск по проиндексированным документам, извлечёт релевантные фрагменты и сгенерирует связный ответ. Если вопрос требует данных из базы, агент сгенерирует SQL-запрос, выполнит его и включит результат в ответ. 

\subsection{Руководство администратора}

\subsubsection{Развёртывание системы}

Чтобы полностью развернуть систему, достаточно из корневой папки проекта выполнить всего одну команду:

\begin{verbatim}
docker compose up -d
\end{verbatim}

После этого станут доступны несколько сервисов. Веб-интерфейс пользователя откроется по адресу \texttt{http://localhost:3000}, основной API RAG Service — по адресу \texttt{http://localhost:8003}. Для управления файлами через консоль MinIO используйте \texttt{http://localhost:9001}, а векторная база ChromaDB будет ждать запросов на \texttt{http://localhost:8000}.

\subsubsection{Настройка параметров индексации}

Основные параметры обработки документов задаются через переменные окружения в файле \texttt{.env}. Например, \texttt{CHUNK\_SIZE=1000} определяет размер фрагмента в символах, \texttt{CHUNK\_OVERLAP=200} — перекрытие между соседними фрагментами, а \texttt{EMBEDDING\_MODEL=intfloat/multilingual-e5-base} указывает, какая модель эмбеддингов используется.

\subsubsection{Мониторинг очереди задач}

Узнать статус фоновой задачи обработки документа можно через специальный эндпоинт: \texttt{GET /status/\{job\_id\}}.

\section{Порядок внедрения и область применения}

\subsection{Целевая аудитория и бизнес-контекст}

Эта система подойдёт любой компании, где есть хранилища данных и возникает потребность быстро получать отчёты по информации из базы данных. Неважно, насколько сложно устроены таблицы и как много в них записей — сотрудники смогут задавать вопросы обычным языком и мгновенно получать готовые цифры и аналитику. Например, руководитель может спросить: «Какие три товара принесли больше всего выручки за прошлый месяц?» — и система сама сформирует и выполнит нужный SQL-запрос, а затем вернёт понятный ответ. Это избавляет аналитиков и менеджеров от необходимости писать сложные запросы или ждать помощи от IT-отдела. Поэтому система будет полезна в розничной торговле, логистике, финансовом секторе, производстве, телекоме и любых других сферах, где данные лежат в базах, а решения нужно принимать быстро.

\section{Характеристика достигнутых результатов}

\subsection{Скорость обработки запросов}

При тестировании мы получили следующие показатели. Векторизация одного запроса занимает от 80 до 150 миллисекунд. Семантический поиск в ChromaDB среди 10 тысяч чанков — от 50 до 100 миллисекунд. Генерация ответа большой языковой моделью (с учётом контекста) и механизма самокоррекции требует от 30 секунд до нескольких минут секунд. Полному агенту, если нужно 1–3 итерации с SQL, требуется 15–45 секунд. Индексация документа на 10 страниц занимает 20–40 секунд.

\subsection{Точность поиска}

Мы используем модель \texttt{intfloat/multilingual-e5-base} с разными префиксами: \texttt{passage:} для документов и \texttt{query:} для запросов. Это позволяет правильно сопоставлять запросы и документы в едином векторном пространстве. Задавая порог косинусного сходства, мы отсеиваем нерелевантные куски и возвращаем только то, что действительно имеет смысл для контекста. А перекрытие чанков в 200 символов помогает не терять смысловые связи на границах фрагментов — так ответ получается более полным и связным.

\subsection{Надёжность агентного модуля}

Агент умеет итеративно улучшать сгенерированный SQL: он делает до пяти попыток, анализирует ошибки и исправляет их самостоятельно. Это обеспечивает устойчивую работу даже при неоднозначных запросах. Кроме того, в DB Connector Service встроена многоуровневая защита: фильтрация опасных ключевых слов, разрешены только \texttt{SELECT}-запросы, установлен таймаут 15 секунд. Всё это исключает риск случайно или намеренно повредить базу данных.

\subsection{Масштабируемость}

Микросервисная архитектура позволяет легко масштабировать самые нагруженные компоненты по горизонтали. Если объём документов растёт, можно увеличить число воркеров, занимающихся индексацией, не меняя архитектуру. Если одновременно работает много пользователей — масштабируем RAG Service и Retrieval Pipeline Service независимо друг от друга. Объём индексируемых данных ограничен только ресурсами сервера: ChromaDB справляется с коллекциями от десятков тысяч до миллионов векторов.

\subsection{Конфиденциальность данных}

Все вычисления происходят локально, внутри корпоративного контура. Модели эмбеддингов запускаются на вашем собственном сервере, файлы хранятся в MinIO, а векторы — в ChromaDB. Корпоративные документы никогда не покидают вашу инфраструктуру, не передаются во внешние облачные сервисы. Это соответствует самым строгим требованиям информационной безопасности большинства организаций.